\chapter{黎曼显式公式}\label{ch:riemann_explicit_formula}

\section{引言：从 Perron 积分到显式公式}\label{sec:riemann_explicit_intro}

经过前文章节的复变分析、积分变换、整函数无穷乘积，以及第~\ref{ch:prime_number_theorem}~章素数定理证明提要，
本章推导黎曼显式公式（Explicit Formula），它将 Chebyshev 阶梯函数 $\psi$ 精确表示为复平面零点级数。

\paragraph{历史位置}
本书的课题始终是素数分布。
第~\ref{ch:elementary_number_theoretic_functions}~章中的 Legendre 型 $x/\log x$ 与 Gauss $\Li(x)$
已是该课题的古典部分；第~\ref{ch:prime_number_theorem}~章证明了主项合法性 $\psi\sim x$（从而 $\pi\sim\Li$）。
黎曼 1859 年的路线（现代严格形式即本章）在同一课题上再进一步，用一条恒等式同时写出：
\begin{itemize}
  \item \textbf{主项}——历史上 Gauss / Legendre 所逼近、并由素数定理保证的那一部分（在 $\psi$ 语言下为 $x$，在 $J$ 语言下为 $\li(x)$）；
  \item \textbf{零点级数}——把 $\pi-\Li$（或 $\psi-x$）的振荡归因于全部非平凡零点 $\{\rho\}$。
\end{itemize}
因而显式公式并非另起炉灶，而是在主项已合法的前提下给出「主项 + 误差结构」的精确表述。
反过来：即使写出该恒等式，也\textbf{不能}单凭公式推出 $\psi\sim x$；
$\sum_{\rho}x^{\rho}/\rho=o(x)$ 依赖于零点实部与可和估计，正是上一章所处理的内容。
恒等与渐近是两层结论，不可互相替代。
黎曼猜想（RH）进一步约束误差阶（von Koch）。本章给出恒等式本身。

在跳跃点取左右平均 $\psi_0$ 后，有
\[
\psi_0(x) = x - \sum_{\rho}\frac{x^{\rho}}{\rho} - \log(2\pi) - \frac12\log(1-x^{-2})
\]
（$x$ 非素数幂时 $\psi_0=\psi$）。左侧是算术阶梯，右侧是遍历 $\zeta(s)$ \textbf{所有非平凡零点}的解析级数，两者通过这一等式建立起精确对应。

公式中，首项 $x$ 给出素数分布的渐近主项（即素数定理的主项，亦是 Legendre--Gauss 主项在 $\psi$ 侧的对应物）；
零点求和项 $\sum x^\rho/\rho$ 则提供围绕主项的振荡修正。
每一个非平凡零点 $\rho$ 在素数阶梯的边缘贡献一个振荡项。
随着求和中纳入零点个数的增加，这些振荡通过叠加逐步逼近阶梯的直角，最终在每个素数处重构出台阶的跳跃。

本章完整推导这一公式：利用 Perron 积分将 Dirichlet 级数化为围道积分，
利用 Hadamard 乘积展开 $\zeta'(s)/\zeta(s)$，并借助留数定理在极点与零点上完成求和。

\section*{本章推导路线图}

\begin{figure}[htbp]
\centering
\resizebox{0.96\textwidth}{!}{%
\begin{tikzpicture}[
  stepbox/.style={draw=MainBlue!70,fill=LightBlue,rounded corners=5pt,
    minimum width=5.0cm,minimum height=1.0cm,
    align=center,font=\small,thick},
  key/.style={draw=Gold!80,fill=LightGold,rounded corners=5pt,
    minimum width=5.0cm,minimum height=1.0cm,
    align=center,font=\small,thick},
  side/.style={draw=DarkGreen!70,fill=LightGreen,rounded corners=4pt,
    minimum width=4.0cm,minimum height=0.95cm,
    align=center,font=\footnotesize},
  light/.style={draw=Gray!60,fill=white,rounded corners=4pt,
    minimum width=5.0cm,minimum height=0.85cm,
    align=center,font=\small,thick},
  arr/.style={-{Stealth[length=6pt]},thick,MainBlue},
  sarr/.style={-{Stealth[length=5pt]},DarkGreen,dashed},
]
% —— 主链 ——
\node[stepbox] (S1) at (0,0)
  {$\zeta(s)$ Euler 乘积\\取对数微分};
\node[stepbox] (S2) at (0,-1.75)
  {$ -\zeta'/\zeta=\sum\Lambda(n)/n^s$\\Dirichlet 级数};
\node[stepbox] (S3) at (0,-3.5)
  {Perron 截断\\$\psi_T(x)=\frac{1}{2\pi i}\int_{c-iT}^{c+iT}
  \left(-\frac{\zeta'}{\zeta}\right)\frac{x^s}{s}\,\mathrm{d}s$};
\node[stepbox] (S4) at (0,-5.25)
  {矩形围道：路径向左平移};
\node[key] (S5) at (0,-7.1)
  {留数定理（计重数）\\$\psi_T=x-\sum_{|\Im\rho|<T}\frac{x^\rho}{\rho}
  -\log(2\pi)+\sum_{n\le N}\frac{x^{-2n}}{2n}+\cdots$};
\node[stepbox] (S6) at (0,-9.0)
  {路径估计：先 $N\to\infty$（$I_3\to0$，平凡零点收全）\\再 $T\to\infty$（$I_2,I_4\to0$）};
\node[key] (S7) at (0,-10.9)
  {$\psi_0(x)=x-\sumrho\frac{x^\rho}{\rho}
  -\log(2\pi)-\frac{1}{2}\log(1-x^{-2})$};
\node[stepbox] (S8) at (0,-12.7)
  {Stieltjes 分部：由 $\psi_0$ 导出 $J_0$};
\node[key] (S9) at (0,-14.5)
  {$ J_0(x)=\li(x)-\sumrho\li(x^\rho)
  -\log 2+\int_x^\infty\frac{\mathrm{d}t}{t(t^2-1)\log t}$};
% 与 $J_0$ 框疏宕拉开，仍以完整实线箭头衔接
\node[light] (S10) at (0,-17.4)
  {Möbius 反演 \eqref{eq:pi_from_J}\\$\pi_0(x)=\sum_{n=1}^{N(x)}\frac{\mu(n)}{n}J_0(x^{1/n})$};

% —— 右侧侧注（均锚向留数/关键公式）——
\node[side] (R1) at (7.0,-5.25)
  {$-\zeta'/\zeta$ 的极点来自\\$\zeta$ 的极点与零点\\{\scriptsize（Hadamard / 第~\ref{ch:xi_function}~章）}};
\node[side] (R2) at (7.0,-7.1)
  {$s=1\to x$\\$\rho\to -m_\rho x^\rho/\rho$\\$-2k\to x^{-2k}/(2k)$\\$s=0\to -\log(2\pi)$};

\foreach \s/\t in {S1/S2,S2/S3,S3/S4,S4/S5,S5/S6,S6/S7,S7/S8,S8/S9,S9/S10}
  \draw[arr] (\s)--(\t);
\draw[sarr] (R1.west)--(S5.east);
\draw[sarr] (R2.west)--(S5.east);
\end{tikzpicture}%
}
\caption{本章显式公式推导主线（至 $J_0$，再经 Möbius 得 $\pi_0$）。
金色框为关键公式；白底框 $\pi_0$ 由第~\ref{ch:elementary_number_theoretic_functions}~章 \eqref{eq:pi_from_J} 代入 $J_0$ 得到（求和至 $N(x)$），不在主留数论证之内。}
\label{fig:explicit-roadmap}
\end{figure}

\section{从 Dirichlet 级数到围道积分}\label{sec:dirichlet-to-contour}

回忆第四章，$\zeta(s)$ 的 Euler 乘积为 $\zeta(s)=\prod_p(1-p^{-s})^{-1}$。取对数微分：
\[
-\frac{\zeta'(s)}{\zeta(s)} = \sum_p \frac{\log p}{p^s-1} = \sum_{n=1}^{\infty} \frac{\Lambda(n)}{n^s},
\]
其中 $\Lambda(n)$ 是 von Mangoldt 函数。这正是第~\ref{ch:riemann_zeta_intro}~章由 Euler 乘积得到的关系。

第~\ref{ch:elementary_number_theoretic_functions}~章定义了 $\psi(x)=\sum_{n\le x}\Lambda(n)$；
其 Dirichlet 级数在 $\Re(s)>1$ 内绝对收敛（第~\ref{ch:riemann_zeta_intro}~章）：
\[
-\frac{\zeta'(s)}{\zeta(s)} = \sum_{n=1}^{\infty} \frac{\Lambda(n)}{n^s}, \qquad \Re(s)>1.
\]

在第\ref{ch:integral_transforms}章中，我们学习了 Perron 公式，它将离散求和转化为围道积分：
\[
\sum_{n\le x} a_n = \frac{1}{2\pi i} \int_{c-i\infty}^{c+i\infty} \left( \sum_{n=1}^{\infty} a_n n^{-s} \right) \frac{x^s}{s} \, ds, \qquad c > \sigma_a.
\]

\begin{lemma}{Perron 核的基本性质}{}
\label{lem:perron_kernel}
对 $c > 0$，Perron 核满足：
\[
  \frac{1}{2\pi i}\int_{c-i\infty}^{c+i\infty}\frac{x^s}{s}\,ds
  = \begin{cases} 1 & x > 1, \\ \tfrac{1}{2} & x = 1, \\ 0 & 0 < x < 1. \end{cases}
\]
直观地说，当 $x>1$ 时，$|x^s|=x^\sigma \to 0$ 随 $\sigma\to-\infty$，故路径向\textbf{左}封闭，包围 $s=0$，$x^s/s$ 在 $s=0$ 处的留数为 $x^0=1$；
当 $x<1$ 时，$|x^s|=x^\sigma \to 0$ 随 $\sigma\to+\infty$，路径向\textbf{右}封闭后无极点，积分为零。
这正是 Perron 公式将 $\sum_{n\le x} a_n$（离散计数）转化为复平面积分的核心机制。
\end{lemma}

在第\ref{ch:riemann_zeta_continuation}章中我们已将 $\zeta(s)$ 解析延拓到全平面（除 $s=1$ 外，见第\ref{ch:riemann_zeta_continuation}章相关定理），在第\ref{ch:zero_distribution}章中我们了解了非平凡零点的分布。在以上的知识铺垫完成后，现在我们可以将 Perron 公式应用于 $\psi(x)$，并对积分线进行围道变形，从而推导出它的显式公式。

如图 \ref{fig:contour} 所示，积分路径从 $\Re(s)=c$（$c>1$）出发，向左平移至 $\Re(s)=-R$，
形成一个矩形围道。当 $R,T\to\infty$ 时，只有右侧垂直线上的积分贡献非零值，
围道内包含 $\zeta(s)$ 的所有极点与非平凡零点。

\begin{figure}[htbp]
\centering
\begin{tikzpicture}[
    scale=0.75,
    >={Stealth[length=6pt,width=5pt]},
    every node/.style={font=\footnotesize}
]
% 临界带填充（浅绿色半透明）
\fill[green!15, draw=none] (0,-6) rectangle (1,6);
% ========== 坐标轴 ==========
\draw[->, thick] (-7,0) -- (3.5,0) node[right] {$\Re(s)$};
\draw[->, thick] (0,-6) -- (0,6) node[above] {$\Im(s)$};

% ========== 坐标轴标签 ==========
\node[below left] at (0,0) {$0$};
\node[below] at (1.7,-0.3) {$c$};
\node[below] at (-4.5,-1.3) {$\sigma_0 = -(2N+1)$};

% ========== c 的位置（c > 1） ==========
\def\realC{1.5} % c = 1.5 > 1
\def\realR{-6.5} % R = 6.5

% ========== 虚部避零高度刻度 ==========
\draw[thick] (-0.1, 3.8) -- (0.1, 3.8);
\node[left, font=\footnotesize, xshift=-2pt] at (0, 3.3) {$iT_0$};
\draw[thick] (-0.1, -3.8) -- (0.1, -3.8);
\node[left, font=\footnotesize, xshift=-2pt] at (0, -3.3) {$-iT_0$};

% ========== 临界线 Re = 1/2 ==========
\draw[thick, brown] (0.5,-6) -- (0.5,6);
\node[brown, right] at (1.0,5.7) {$\Re(s)=\frac12$（临界线）};

% ========== 原始积分路径（红色实线，画在 c = 1.5） ==========
\draw[solid, red!70!black, thick] (\realC,-4.5) -- (\realC,4.5);
\node[red!70!black, right] at (\realC,4.2) {$\Re(s)=c\;(c>1)$（原始积分路径）};

% ========== 矩形围道（右边线与红色虚线重合） ==========
\coordinate (A) at (\realC, -3.8);
\coordinate (B) at (\realC, 3.8);
\coordinate (C) at (\realR, 3.8);
\coordinate (D) at (\realR,-3.8);

\draw[thick, blue!70!black, 
      decoration={markings, mark=at position 0.18 with {\arrow{>}},
                         mark=at position 0.42 with {\arrow{>}},
                         mark=at position 0.68 with {\arrow{>}},
                         mark=at position 0.92 with {\arrow{>}}},
      postaction={decorate}]
    (A) -- (B) -- (C) -- (D) -- cycle;

% ========== 围道各边标注 ==========
\node[blue!70!black, above] at (-3.0, 4.0) {$\Gamma_2$（上水平边 $\Im(s)=T_0$）};
\node[blue!70!black, below] at (-3.2,-4.0) {$\Gamma_4$（下水平边 $\Im(s)=-T_0$）};
\node[blue!70!black, left] at (-3.0, 1.2) {$\Gamma_3$（左垂直边）};
\node[blue!70!black, right] at (1.6,-1.3) {$\Gamma_1$（右垂直边）};

% ========== 极点标记 ==========
% 非平凡零点（示意画于临界线 Re=1/2 上；实际仅知位于临界带 0<Re(s)<1 内）
\foreach \y in {1.2, 2.3, 3.1, -1.1, -2.0, -3.3} {
    \filldraw[red] (0.5,\y) circle (2.5pt);
}
\node[red, right] at (1.5,3.3) {$\rho$（非平凡零点，示意）};
\node[red, right] at (1.5,-3.3) {$\bar\rho$（共轭零点，示意）};

% s=1 处的极点
\filldraw[orange!80!black] (1,0) circle (3pt);
\node[orange!80!black, above right] at (1.5,0.15) {$s=1$（极点）};

% 平凡零点（左半平面）
\foreach \n in {1,2,3} {
    \filldraw[blue!50!black] (-2*\n,0) circle (2pt);
}
\node[blue!50!black, below] at (-2,-0.1) {$-2$};
\node[blue!50!black, below] at (-4,-0.1) {$-4$};
\node[blue!50!black, below] at (-6,-0.1) {$-6$};
\node[blue!50!black] at (-3.5,-1.0) {平凡零点};

% s=0 点
\filldraw[blue!50!black] (0,0) circle (2pt);
\node[blue!50!black, above left] at (0,0.15) {$s=0$};

% ========== 围道内提示 ==========
\node[gray!60!black, font=\footnotesize] at (-2.5,-2.5) {围道内含所有极点};

% ========== 水平边积分趋零的说明 ==========
\node[gray!60!black, font=\footnotesize] at (-3, 3.0) {$T_0\to\infty$ 时};
\node[gray!60!black, font=\footnotesize] at (-3, 2.5) {$\Gamma_2,\Gamma_4\to 0$};

% ========== 图例（加框；行距压缩以降低框高） ==========
\begin{scope}[shift={(6.8, -0.5)}]
    \draw[thick, gray!70!black, fill=white, rounded corners=4pt]
        (-1.7, 3.45) rectangle (6.2, -1.35);

    \draw[thick, blue!70!black,
          decoration={markings, mark=at position 0.5 with {\arrow{>}}},
          postaction={decorate}]
        (-1.3, 2.95) -- (-0.3, 2.95);
    \node[blue!70!black, right, font=\footnotesize] at (-0.1, 2.95) {围道方向（逆时针）};

    \draw[solid, red!70!black, thick] (-1.3, 2.35) -- (-0.3, 2.35);
    \node[red!70!black, right, font=\footnotesize] at (0.1, 2.35) {原始路径 $\Re(s)=c\;(c>1)$};

    \filldraw[orange!80!black] (-0.8, 1.75) circle (2.5pt);
    \node[orange!80!black, right, font=\footnotesize] at (0, 1.75) {$s=1$ 极点（主项 $x$）};

    \filldraw[red] (-0.8, 1.15) circle (2.2pt);
    \node[red, right, font=\footnotesize] at (-0.1, 1.15) {非平凡零点 $\rho$（振荡项）};

    \filldraw[blue!50!black] (-0.8, 0.55) circle (2pt);
    \node[blue!50!black, right, font=\footnotesize] at (-0.1, 0.55) {$s=0,-2,-4,\dots$（修正项）};

    \draw[thick, brown] (-1.3, -0.05) -- (-0.3, -0.05);
    \node[brown, right, font=\footnotesize] at (-0.1, -0.05) {临界线 $\Re(s)=1/2$};

    \fill[green!15, draw=none] (-1.3, -0.85) rectangle (-0.3, -1.15);
    \node[green!70!black, right, font=\footnotesize] at (0.3, -1.0) {临界带 $0<\Re(s)<1$};
\end{scope}

\end{tikzpicture}
\caption{显式公式推导中的矩形围道（$c>1$，右边线位于 $\Re(s)=c$），左边界位于 $\Re(s) = \sigma_0 = -(2N+1)$。围道水平边 $\Gamma_2, \Gamma_4$ 位于避零高度 $\Im(s) = \pm T_0$ 处以保证路径估计收敛。图中非平凡零点 $\rho$ 仅为示意，一般情形下仅知其位于临界带内，推导不依赖 RH。}
\label{fig:contour}
\end{figure}

\section{显式公式的推导}\label{sec:explicit-derivation}

\subsection{应用 Perron 积分公式}

出发点与第~\ref{ch:prime_number_theorem}~章相同：第~\ref{ch:discrete_continuous}~章的~$\psi$ 的 Perron 表示~\eqref{eq:psi_perron}
（$c>1$；非素数幂处写 $\psi(x)$，跳跃点取平均）：
\[
\psi(x)=\frac{1}{2\pi i}\int_{c-i\infty}^{c+i\infty}\left(-\frac{\zeta'(s)}{\zeta(s)}\right)\frac{x^s}{s}\,ds.
\]
下文围道变形即从该式出发；\textbf{不再另编号}，以免与~\eqref{eq:psi_perron} 重复。

\paragraph{Perron 截断}
有限高度版本为
\[
\sum_{n\le x} a_n
=\frac{1}{2\pi i}\int_{c-iT}^{c+iT} F(s)\frac{x^s}{s}\,ds
+O\!\Bigl(\frac{x^c}{T}\sum_n\frac{|a_n|}{n^c}\Bigr)
\qquad(c>\sigma_a).
\]
当 $a_n=\Lambda(n)$、$c>1$ 时截断误差 $O(x^c/T)$；令 $T\to\infty$ 得~\eqref{eq:psi_perron}
（跳跃点取平均，详见第~\ref{ch:integral_transforms}~章）。

\paragraph{$\psi(x)$ 的 Mellin 积分变换}
作为 Perron 积分~\eqref{eq:psi_perron} 的逆向桥梁，第二 Chebyshev 函数 $\psi(x)$ 的 Mellin 积分变换将实轴上的素数计数阶梯还原为复变 $\zeta$ 导数之比。对任意 $\Re(s) > 1$，有：
\begin{equation}
  -\frac{\zeta'(s)}{\zeta(s)} \;=\; s\int_{1}^{\infty}\frac{\psi(x)}{x^{s+1}}\,\mathrm{d}x.
\label{eq:psi_mellin}
\end{equation}
这一关系在第~\ref{ch:spectrum_map}~章图~\ref{fig:function-relations} 上以连线 $\psi(x) \leftrightarrow \zeta(s)$（Mellin）标注，代表数论加权计数与复解析性质之间的对偶桥梁。


\subsection{围道变形}

本节的推导采用\textbf{矩形围道}，便于与竖直直线上的 Perron 积分衔接。
黎曼 1859 年原稿在素数公式一侧采用的是竖直线反演（Fourier / 逆 Mellin 型）论证，
而非用 Hankel 围道推出 $\psi_0$；Hankel 围道主要用于 $\zeta$ 的解析延拓（第~\ref{ch:riemann_zeta_continuation}~章第一法）。
矩形围道与原稿竖直线在极限下包围同一类奇点贡献，故最终显式公式一致。

\textbf{注意}：原始 Perron 积分~\eqref{eq:psi_perron} 的积分限为 $\pm i\infty$，无法直接构成有界闭合矩形。因此，我们首先对有限截断版本 $\psi_T(x)$（见下面第~\ref{sec:finite-truncation} 节，公式 \eqref{eq:psi_T_formula}）进行围道推导，再令 $T\to\infty$。

对截断积分，将路径（$\Re(s)=c$，$|\Im s|\le T$）向左平移至 $\Re(s)=-R$（$R=2N+1$ 选在相邻平凡零点之间），加上上下水平线段和左侧垂直线段，形成闭合矩形围道。对此有界矩形，由留数定理，其积分等于围道内所有极点的留数之和（含 $|\Im\rho|<T$ 的非平凡零点、$s=1$ 处的极点、平凡零点 $-2,-4,\ldots,-2N$ 及 $s=0$）。

令 $N\to\infty$、$T\to\infty$（具体极限顺序见第~\ref{sec:path-estimation} 节），左侧垂直边以及两条水平边上的积分均趋于零（需利用 $\zeta(s)$ 在左半平面的增长估计，详见第~\ref{sec:path-estimation} 节）。因此 $\psi_T(x)$ 在 $T\to\infty$ 后趋于所有极点的留数之和。

\subsection{极点、留数与有限截断}\label{sec:residue-sum}
\label{sec:finite-truncation}

被积函数 $F(s)=\bigl(-\zeta'(s)/\zeta(s)\bigr)\,x^s/s$ 的极点与留数如下
（非平凡零点 $\rho$ 的重数记为 $m_\rho$；下文求和默认\textbf{计重数}，
单零点时 $m_\rho=1$，留数贡献为 $-m_\rho x^\rho/\rho$）。

\paragraph{极点在 $s=1$}
因 $\zeta(s)=\frac{1}{s-1}+\gamma+O(s-1)$，有 $-\zeta'/\zeta=1/(s-1)+O(1)$，故
\[
\operatorname{Res}_{s=1}\!\left[\left(-\frac{\zeta'(s)}{\zeta(s)}\right)\frac{x^s}{s}\right]
=\frac{x}{1}=x.
\]

\paragraph{极点在非平凡零点 $\rho$（重数 $m_\rho$）}
$\zeta(s)\sim c_\rho(s-\rho)^{m_\rho}$ 时，$-\zeta'/\zeta$ 在 $\rho$ 处留数为 $-m_\rho$，从而
\[
\operatorname{Res}_{s=\rho}\!\left[\left(-\frac{\zeta'(s)}{\zeta(s)}\right)\frac{x^s}{s}\right]
= -m_\rho\frac{x^{\rho}}{\rho}.
\]
（下文常将 $m_\rho$ 并入 $\sum_\rho$，写作 $\sum_\rho x^\rho/\rho$ 并声明计重数。）

\paragraph{极点在平凡零点 $s=-2n$（$n=1,2,\ldots$）}
一阶零点时
\[
\operatorname{Res}_{s=-2n}\!\left[\left(-\frac{\zeta'(s)}{\zeta(s)}\right)\frac{x^s}{s}\right]
=\frac{x^{-2n}}{2n}.
\]
对全部平凡零点求和，令 $y=x^{-2}$（$x>1$），得
\[
\sum_{n=1}^{\infty}\frac{x^{-2n}}{2n}
=-\frac12\log(1-x^{-2}).
\]

\paragraph{极点在 $s=0$}
$\zeta(0)=-\frac12\neq 0$，$\zeta'(0)=-\frac12\log(2\pi)$；极点仅来自 $1/s$，
\[
\operatorname{Res}_{s=0}\!\left[\left(-\frac{\zeta'(s)}{\zeta(s)}\right)\frac{x^s}{s}\right]
=-\frac{\zeta'(0)}{\zeta(0)}=-\log(2\pi).
\]

\subsection{有限截断与极限过渡}


由于级数 $\sum_{\rho} \frac{x^{\rho}}{\rho}$ 条件收敛（见第~\ref{sec:convergence} 节），我们首先考虑有限截断版本。设
\[
\psi_T(x) = \frac{1}{2\pi i} \int_{c-iT}^{c+iT} \left(-\frac{\zeta'(s)}{\zeta(s)}\right) \frac{x^s}{s} ds.
\]
通过围道变形和留数定理，矩形围道（截断高度 $T$，左边 $\sigma_0=-(2N+1)$）内所含极点贡献了以下有限截断形式：

\begin{equation}
\psi_{T,N}(x) = x - \sum_{|\Im\rho|<T} \frac{x^{\rho}}{\rho} - \log(2\pi) + \sum_{n=1}^{N}\frac{x^{-2n}}{2n} + \text{（三条补充边积分的误差项）}.
\label{eq:psi_T_formula}
\end{equation}

其中 $s=1$ 的留数贡献了 $x$，非平凡零点贡献了 $-\sum_{|\Im\rho|<T}x^\rho/\rho$，$s=0$ 的留数贡献了 $-\log(2\pi)$，平凡零点 $s=-2,-4,\ldots,-2N$ 各贡献 $x^{-2n}/(2n)$（有限个）。

先令 $N\to\infty$（左边 $\sigma_0\to-\infty$）：补充路径估计中左竖线积分趋于零，同时平凡零点级数收全，$\sum_{n=1}^N x^{-2n}/(2n)\to-\tfrac12\log(1-x^{-2})$，故：
\[
\psi_T(x) = x - \sum_{|\Im\rho|<T} \frac{x^{\rho}}{\rho} - \log(2\pi) - \tfrac12\log(1-x^{-2}) + \text{（水平边积分误差项）}.
\]

再令 $T\to\infty$：由下面的补充路径估计（见下一子节），水平边 $I_2, I_4$ 上的积分均趋于零，故
$\psi_T(x)\to\psi_0(x)$（非素数幂处即 $\psi(x)$），极限存在。

\subsection{补充路径估计}\label{sec:path-estimation}

以下三个估计确保矩形围道的三条补充边（顶边 $I_2$、底边 $I_4$、左竖边 $I_3$）上的积分在 $T\to\infty$ 时趋于零，从而使有限截断 $\psi_T(x)$ 极限过渡到完整的显式公式成立。

为了严格证明式 \eqref{eq:psi_T_formula} 中补充路径积分趋于零，
需要以下三个估计。


\paragraph{估计 $|\zeta'/\zeta(\sigma+iT)| \ll (\log T)^2$ 的依据}
在复变函数的估计中，为了得到对数导数在临界带内的水平增长上界，直接应用一般的复估计（如 Phragmén--Lindelöf 原理）是困难的，因为零点的存在会导致 $\zeta'/\zeta$ 在奇点处发散。为此，本书采用逻辑闭环的\textbf{“零点避让空隙法”}，其推导直接基于前述两个定理：

\begin{enumerate}
  \item \textbf{Hadamard 部分分式展开}（参见第\ref{ch:xi_function}章关于 Hadamard 乘积的讨论）：
        在临界带 $-1 \le \sigma \le 2$ 内部，$\zeta(s)$ 的对数导数具有如下精确的部分分式表示：
        \[
        \frac{\zeta'(s)}{\zeta(s)} = \sum_{|t - \gamma| \le 1} \frac{1}{s - \rho} + O(\log |t|)
        \]
        其中求和项跑遍所有虚部满足 $|t-\gamma| \le 1$ 的非平凡零点 $\rho = \beta + i\gamma$。该展开式表明，对数导数的模主要取决于其与邻近零点的距离。
  \item \textbf{Riemann--von Mangoldt 零点计数定理与鸽巢原理}：
        在高度区间 $[T, T+1]$ 内部，零点的总个数约为 $O(\log T)$。根据鸽巢原理，若将此区间平均分为若干个空隙，则必然存在一个高度 $T_0 \in [T, T+1]$，使得它到任何零点虚部 $\gamma$ 的垂直距离至少为：
        \[
        |T_0 - \gamma| \gg \frac{1}{\log T}
        \]
        如果我们在此“避零高度” $t = T_0$ 上估计部分分式展开式中的分母，则有下界估计 $|s - \rho| \ge |T_0 - \gamma| \gg 1/\log T$。由此，对展开式中的每一项有 $\frac{1}{|s-\rho|} \ll \log T$。
\end{enumerate}

由于在高度满足 $|T_0 - \gamma| \le 1$ 的求和项共有 $O(\log T)$ 个，我们将它们直接累加，便有：
\[
\left| \sum_{|T_0 - \gamma| \le 1} \frac{1}{s - \rho} \right| \le \sum_{|T_0 - \gamma| \le 1} \frac{1}{|s-\rho|} \ll \log T \cdot \log T = (\log T)^2
\]
结合余项 $O(\log T)$，最终可得对 $\sigma \ge -1$ 且选取此避零高度 $T_0$ 时，对数导数的上界满足 $\left|\frac{\zeta'}{\zeta}(\sigma+iT_0)\right| = O\bigl((\log T_0)^2\bigr)$。

由于该避零高度可在任意大的 $T$ 附近微调选取，后续将 $T$ 取为避零高度，在 $\sigma\ge -1$ 的水平边上有
\[
\left|\frac{\zeta'}{\zeta}(\sigma+iT)\right| = O\bigl((\log T)^2\bigr)
\]
（标准上界的\textbf{提要}；完整证明见 Ingham / Titchmarsh 等专著）。

\subsection*{（1）顶边 $I_2$ 与底边 $I_4$ 的估计}

以顶边 $I_2$ 为例（底边 $I_4$ 类似）。沿 $s=\sigma+iT$，
$\sigma$ 从 $c$ 到 $\sigma_0$（$\sigma_0<0$）：
\[
I_2 = \int_c^{\sigma_0} \left(-\frac{\zeta'}{\zeta}(\sigma+iT)\right) 
      \frac{x^{\sigma+iT}}{\sigma+iT}\, d\sigma.
\]

由无零点区域定理（$\zeta(s)$ 在 $\sigma\ge 1-c_1/\log T$ 内无零点），
有估计：
\[
\left|\frac{\zeta'}{\zeta}(\sigma+iT)\right| \ll (\log T)^2,
\quad \sigma\ge -1,\ |t|=T.
\]
代入得：
\[
|I_2| \ll \int_{\sigma_0}^{c} (\log T)^2 \cdot \frac{x^\sigma}{T}\, d\sigma
\ll \frac{(\log T)^2}{T} \cdot \frac{x^c - x^{\sigma_0}}{\log x}
\xrightarrow{T\to\infty} 0.
\]

\subsection*{（2）左竖线 $I_3$ 的估计}

沿 $s=\sigma_0+it$，$t$ 从 $-T$ 到 $T$（$\sigma_0=-(2N+1)$）：
\[
I_3 = \int_{-T}^{T} \left(-\frac{\zeta'}{\zeta}(\sigma_0+it)\right) 
      \frac{x^{\sigma_0+it}}{\sigma_0+it}\, dt.
\]

由泛函方程与凸性估计，在 $\sigma_0=-(2N+1)$ 的左竖线上有
\[
\left|\frac{\zeta'}{\zeta}(\sigma_0+it)\right| \ll \log |t|,\quad |t|\ge 2
\]
（提要级上界）。
于是：
\[
|I_3| \ll x^{\sigma_0} \int_{-T}^{T} \frac{\log(|t|+2)}{|\sigma_0+it|}\, dt
\ll x^{\sigma_0} \cdot T \log T
\xrightarrow{\sigma_0\to-\infty} 0
\]
（对固定 $T$，先令 $N\to\infty$ 使 $\sigma_0\to -\infty$）。


\paragraph{两步极限的顺序}
三条补充边的估计对应两个独立的极限参数，其正确操作顺序如下：
\begin{enumerate}
  \item \textbf{先固定 $T$，令 $N\to\infty$（$\sigma_0\to-\infty$）}：
        左竖线 $I_3\to 0$，同时所有平凡零点 $s=-2,-4,\dots,-2N$ 的留数被收全，
        其尾项 $\sum_{k>N} x^{-2k}/(2k)\to 0$（由等比级数估计）。
  \item \textbf{再令 $T\to\infty$}：
        上下水平边 $I_2, I_4\to 0$（由估计（1）），
        同时所有虚部 $|\Im\rho|<T$ 的非平凡零点留数被逐步收入，
        其部分和在对称求和意义下收敛（由第\ref{ch:xi_function}章第\ref{sec:xi-hadamard-product}节关于零点对称配对与条件收敛的讨论）。
\end{enumerate}
两步极限完成后，$\psi_T(x)\to\psi_0(x)$，即有限截断公式 \eqref{eq:psi_T_formula} 趋于显式公式 \eqref{eq:explicit_psi}。


\subsection{综合：显式公式}

将以上所有留数相加，并令 $R\to\infty$，得到：

\begin{theorem}{Riemann--von Mangoldt 显式公式}{}
对 $x>1$，令 $\psi_0(x)=\frac{\psi(x^-)+\psi(x^+)}{2}$（非素数幂处 $\psi_0=\psi$）。则
\begin{equation}
\boxed{\psi_0(x)=x-\sum_{\rho}\frac{x^{\rho}}{\rho}-\log(2\pi)-\frac12\log(1-x^{-2})}
\label{eq:explicit_psi}
\end{equation}
其中 $\sum_{\rho}$ 遍历 $\zeta$ 的全部非平凡零点，\textbf{计重数}，
并按 $|\Im\rho|$ 对称截断求和：
$\sum_{\rho}=\lim_{T\to\infty}\sum_{|\Im\rho|<T}$（同一 $\rho$ 出现 $m_\rho$ 次，或等价地写 $\sum_\rho m_\rho x^\rho/\rho$）。
有限截断 $\psi_T$ 的路径误差在 $T\to\infty$ 时趋于零（第~\ref{sec:path-estimation}~节，估计为提要级）。
（跳跃点 $\psi_0$ 见书前《符号与约定》；$\Li$ 与 $\li$ 的分工见第~\ref{ch:logint_Ei}~章。）
\end{theorem}


\paragraph{关于跳跃点与 $\psi/\psi_0$}
当 $x$ 为素数幂 $p^k$ 时，$\psi$ 有跳跃，\eqref{eq:explicit_psi} 给出的是左右平均 $\psi_0(x)$。
叙述中若声明“$x$ 非素数幂”，可写 $\psi(x)$ 代替 $\psi_0(x)$。


\section{分部求和公式的显式推导：从 \texorpdfstring{$\psi(x)$}{psi(x)} 到 \texorpdfstring{$J(x)$}{J(x)}}\label{sec:integration-by-parts}

阶梯函数上的 Stieltjes 积分与分部公式的\textbf{一般表述}见附录~\ref{app:abel_stieltjes}
（定义~\ref{def:app-stieltjes}、定理~\ref{thm:app-stieltjes-parts}）。
本节将一般公式专用于 $\psi\to J$。

回顾 $J(x)$ 与 $\psi(x)$ 的定义
（第~\ref{ch:elementary_number_theoretic_functions}~章；$J$ 即黎曼原稿中的 $f$）：
\[
\psi(x) = \sum_{p^k\le x} \log p,\qquad 
J(x) = \sum_{p^k\le x} \frac{1}{k}.
\]

注意到 $J(x)$ 也可写作：
\[
J(x) = \sum_{n\le x} \frac{\Lambda(n)}{\log n},
\]
因为当 $n=p^k$ 时 $\Lambda(n)=\log p$，$\log n = k\log p$，其比值恰为 $1/k$；而当 $n$ 非素数幂时 $\Lambda(n)=0$，故对应项贡献为零；$n=1$ 时 $\Lambda(1)=0$，故该项亦无贡献，可直接略去。因此求和虽形式上遍历所有 $n\le x$，实际上只有素数幂项有非零贡献，与 $J(x)=\sum_{p^k\le x} 1/k$ 的定义完全等价。

\subsection*{化为跳跃和}

由附录~\ref{app:abel_stieltjes}，对阶梯 $\psi$ 有
\[
\int_{2^-}^x f(t)\, \mathrm{d}\psi(t) = \sum_{2 \le n \le x} f(n)\Lambda(n).
\]
取 $f(t)=1/\log t$ 即
\[
\int_{2^-}^{x} \frac{1}{\log t}\, \mathrm{d}\psi(t)
= \sum_{p^k \le x} \frac{\log p}{\log(p^k)}
= \sum_{p^k \le x} \frac{1}{k} = J(x),
\]
与第~\ref{ch:elementary_number_theoretic_functions}~章 \eqref{eq:J_from_psi_stieltjes} 一致。

\subsection*{分部积分}

附录~\ref{app:abel_stieltjes} 式~\eqref{eq:app-stieltjes-parts}：若 $f\in C^1$，$A$ 右连续阶梯，则
\[
\int_a^b f(t)\, \mathrm{d}A(t) = f(b)A(b) - f(a)A(a^-) - \int_a^b A(t)f'(t)\, \mathrm{d}t.
\]
取 $a=2$，$b=x$，$f(t)=1/\log t$，$A(t)=\psi(t)$，则
\[
J(x) = \frac{\psi(x)}{\log x} - \frac{\psi(2^-)}{\log 2} - \int_2^x \psi(t)\cdot \frac{d}{dt}\!\left(\frac{1}{\log t}\right) dt.
\]
其中 $\psi(2^-)=\lim_{t\to 2^-}\psi(t)$。由于 $\psi(t)=0$ 对所有 $t<2$ 成立，故 $\psi(2^-)=0$。

\subsection*{计算微分项}

\[
\frac{d}{dt}\!\left(\frac{1}{\log t}\right) = -\frac{1}{t(\log t)^2}.
\]

代入得：
\[
J(x) = \frac{\psi(x)}{\log x} + \int_2^x \frac{\psi(t)}{t(\log t)^2}\, dt.
\]

由于 $\psi(2^-)=0$，最终：
\begin{equation}
\boxed{J(x) = \frac{\psi(x)}{\log x} + \int_2^x \frac{\psi(t)}{t(\log t)^2}\, dt.}
\label{eq:J_from_psi_explicit}
\end{equation}

\subsection*{代入 $\psi(x)$ 的显式公式}

将 $\psi$ 的显式公式 \eqref{eq:explicit_psi}（在 $x$ 非素数幂处即 $\psi=\psi_0$）
\[
\psi(x)=x-\sum_\rho \frac{x^\rho}{\rho}-\log(2\pi)-\frac12\log(1-x^{-2})
\]
代入 \eqref{eq:J_from_psi_explicit}，逐项计算。

\paragraph{（1）主项 \texorpdfstring{$t$}{t} 的贡献}
利用恒等式 $ \int_2^x \frac{dt}{(\log t)^2} = \li(x) - \li(2) - \frac{x}{\log x} + \frac{2}{\log 2}$，其中 $\li(x)=\int_0^x \frac{dt}{\log t}$，可得：
\[
\begin{aligned}
\frac{x}{\log x} + \int_2^x \frac{t}{t(\log t)^2}\, dt
&= \frac{x}{\log x} + \int_2^x \frac{dt}{(\log t)^2} \\
&= \frac{x}{\log x} + \left( \li(x) - \li(2) - \frac{x}{\log x} + \frac{2}{\log 2} \right) \\
&= \li(x) - \li(2) + \frac{2}{\log 2}.
\end{aligned}
\]

\paragraph{（2）常数项 \texorpdfstring{$-\log(2\pi)$}{-log(2 pi)} 的贡献}
利用 $ \int_2^x \frac{dt}{t(\log t)^2} = \frac{1}{\log 2} - \frac{1}{\log x}$，得：
\[
-\frac{\log(2\pi)}{\log x} - \log(2\pi) \int_2^x \frac{dt}{t(\log t)^2}
= -\frac{\log(2\pi)}{\log x} - \log(2\pi)\left( \frac{1}{\log 2} - \frac{1}{\log x} \right)
= -\frac{\log(2\pi)}{\log 2}.
\]

\paragraph{（3）平凡零点项 \texorpdfstring{$g(t)=-\frac12\log(1-t^{-2})$}{g(t) = -1/2 log(1-t^-2)} 的贡献}
将 $g(t)$ 展开为级数 $g(t)=\sum_{k=1}^\infty \frac{t^{-2k}}{2k}$，利用与主项相同的积分恒等式，逐项求和可得（两端相差与 $x$ 无关的积分下限常数，该常数将在后文统一合并）：
\[
\frac{g(x)}{\log x} + \int_2^x \frac{g(t)}{t(\log t)^2}\, dt
= \int_x^\infty \frac{dt}{t(t^2-1)\log t} + \text{常数}.
\]

\paragraph{（4）非平凡零点项 \texorpdfstring{$-\sum_\rho t^\rho/\rho$}{-sum t^rho / rho} 的贡献}
对固定 $\rho$（$0<\operatorname{Re}\rho<1$），记 $\operatorname{li}(x^\rho)$ 为沿适当路径的解析延拓
（等价于 $\operatorname{Ei}(\rho\log x)$ 的主支选取，见 Edwards）。
关键恒等式为：当 $x>1$ 时，
\[
\frac{\mathrm{d}}{\mathrm{d}x}\operatorname{li}(x^\rho)
=
\frac{x^{\rho-1}}{\log x}.
\]
对 $J$ 的积分表示作 Stieltjes 分部后，单项
\[
-\frac{x^\rho/\rho}{\log x} - \int_2^x \frac{t^\rho/\rho}{t(\log t)^2}\, dt
\]
与 $-\operatorname{li}(x^\rho)$ 仅差一个与 $x$ 无关的下限常数：
可对 $\operatorname{li}(x^\rho)$ 再作一次分部积分（或直接对参数 $\rho$ 求导并与
$\int_2^x t^{\rho-1}/(\log t)^2\,\mathrm{d}t$ 对照）得到
\[
-\frac{x^\rho/\rho}{\log x} - \int_2^x \frac{t^\rho/\rho}{t(\log t)^2}\, dt
= -\li(x^\rho) + C(\rho),
\]
其中 $C(\rho)$ 由 $t=2$ 处的边界项给出。对所有零点求和，常数项与前述常数合并
（详细数值恒等见 Edwards (1974) §1.15）。


\paragraph{常数项合并说明}
各项贡献产生的与 $x$ 无关的常数，其来源概述如下：
\begin{itemize}
  \item \textbf{主项 $t$ 的常数贡献}：$-\li(2) + \frac{2}{\log 2}$（来自积分下限 $t=2$）。
  \item \textbf{常数项 $-\log(2\pi)$ 的贡献}：$-\frac{\log(2\pi)}{\log 2}$。
  \item \textbf{平凡零点项（尾项积分）的常数贡献}：有限下限 $t=2$ 处产生的常数，以及将
  $\sum_{n=1}^{\infty}\frac{x^{-2n}}{2n}=-\frac12\log(1-x^{-2})$
  经分部积分后改写为
  $\int_x^{\infty}\frac{dt}{t(t^2-1)\log t}$ 时多出的与 $x$ 无关的部分。
  （该正值积分本身\textbf{没有}简单的错误闭式；其与其它边界常数的合并见 Edwards (1974) §1.15，此处不写不可核验的显式数值恒等式。）
  \item \textbf{非平凡零点项的常数贡献}：每个零点 $\rho$ 的积分在下限 $t=2$ 处各贡献一个 $\rho$ 相关常数 $-\li(2^\rho)$，其全体对称求和后（利用 Hadamard 乘积对数导数展开）与 $\li(2)$ 等项一并参与合并（详见 Edwards (1974) §1.15）。
\end{itemize}
将以上各项与 $x$ 无关的常数相加，Edwards 的完整计算表明它们最终合并为 $-\log 2$，得到公式中的最后一项。


\paragraph{（5）常数合并}
将以上所有贡献相加。主项、$-\log(2\pi)$、平凡零点项与各 $\li(2^\rho)$ 等
\textbf{与 $x$ 无关的下限常数}，经完整计算合并为单一常数 $-\log 2$
（Edwards (1974) §1.15；本章给出来源分类，\textbf{不在此重算}该数值恒等）。
于是得到 $J$ 的显式公式（$x$ 非素数幂时 $J_0=J$）：
\begin{equation}
\boxed{J_0(x) = \li(x) - \sum_\rho \li(x^\rho) + \int_x^\infty \frac{dt}{t(t^2-1)\log t} - \log 2,}
\label{eq:explicit_J_final}
\end{equation}
其中 $\li$ 为 Cauchy 主值对数积分（\textbf{非}偏移版 $\Li$；
实变量上 $\Li$ 与 $\li$ 的关系见第~\ref{ch:logint_Ei}~章式~\eqref{eq:Li-li-real}；
$\sum_\rho$ 计重数并对称截断。常数 $-\log 2$ 与 $\li$ 配套，不可在保留该常数的同时把 $\li$ 换成 $\Li$。

\paragraph{显式公式中出现的是哪一种}
本节落实第~\ref{ch:logint_Ei}~章第~\ref{sec:Li-li-convention}~节的约定：
\textbf{素数定理比较用 $\Li$（实），显式公式用 $\li$（可含 $x^{\rho}$）}。
\begin{enumerate}
  \item \textbf{$\psi$ 形}（Riemann--von Mangoldt）：
    主项为 $x$，不出现 $\Li$ 或 $\li$；见式~\eqref{eq:explicit_psi}。
  \item \textbf{$J$ 形}（经典黎曼\,/\,Edwards）：
    主项为 $\li(x)$，零点项为 $\li(x^{\rho})$，并含常数 $-\log 2$ 等；
    见式~\eqref{eq:explicit_J_final}。
  \item \textbf{$\pi$ 形}（Möbius 反演后）：
    主项族为 $\sum_{n}\mu(n)/n\,\li(x^{1/n})$ 等，仍使用 $\li$，
    不使用 $\Li(x^{\rho})$。
\end{enumerate}

\paragraph{$\li(x^\rho)$ 与级数收敛性}
复函数 $\li(z)$ 与 $\Ei$ 见第~\ref{ch:logint_Ei}~章第~\ref{sec:li-complex-power}~节；
此处 $\li(x^{\rho})$ 即在点 $z=x^{\rho}$ 上取值（$x^{\rho}\notin[1,\infty)$ 时由该章定义给出），不另立新函数。
单个 $\li(x^{\rho})$ 一般是复数；
当 $|\Im\rho|\to\infty$ 时，由 $\li(z)\sim z/\log z$ 得
$\li(x^{\rho})\sim x^{\rho}/(\rho\log x)$。

级数 $\sum_{\rho} \li(x^\rho)$ 的收敛性需要特别处理：
\begin{itemize}
  \item 该级数\textbf{不绝对收敛}，必须按 $|\Im(\rho)|$ 递增的顺序对称求和：
        $\sum_{\rho} = \lim_{T\to\infty} \sum_{|\Im\rho|<T}$。
  \item 由 $\li(x^\rho) \sim \frac{x^\rho}{\rho\log x}$（当 $|\Im\rho|\to\infty$），
        结合第\ref{ch:xi_function}章第\ref{sec:xi-hadamard-product}节阐明的 $\sum 1/|\rho|^2$ 绝对收敛性，可知在对称求和意义下该级数条件收敛。
\end{itemize}
这一收敛性与 $\psi(x)$ 显式公式中的 $\sum x^\rho/\rho$ 完全平行。

\paragraph{共轭零点配对后为实}
因 $\zeta$ 的 Dirichlet 系数实，非平凡零点成共轭对出现：
若 $\rho=\beta+i\gamma$ 为零点，则 $\bar\rho=\beta-i\gamma$ 亦为零点。
\begin{enumerate}
  \item \textbf{$\psi$ 形：}
    对 $x>1$ 实，
    \[
      \frac{x^{\rho}}{\rho}+\frac{x^{\bar\rho}}{\bar\rho}
      \;=\;
      2\operatorname{Re}\!\left(\frac{x^{\rho}}{\rho}\right)
      \;\in\;
      \mathbb{R}.
    \]
  \item \textbf{$J$ 形：}
    在 $\li$ 与 $\Ei$ 的主支与共轭相容时，
    \[
      \li(x^{\rho})+\li(x^{\bar\rho})
      \;=\;
      2\operatorname{Re}\li(x^{\rho})
      \;\in\;
      \mathbb{R}.
    \]
\end{enumerate}
因此：\textbf{单个} $\li(x^{\rho})$ 或 $x^{\rho}/\rho$ 一般是复数；
\textbf{共轭配对}（并对称截断）后贡献为实，与 $\psi_0(x)$、$J_0(x)$ 为实一致。


\section{从 \texorpdfstring{$J(x)$}{J(x)} 到 \texorpdfstring{$\pi(x)$}{pi(x)}}\label{sec:psi-to-pi}

本节考虑 $x>2$ 的情形，此时 $\li(x)$ 和被积函数均无奇点问题。

显式公式不仅适用于 $\psi(x)$，也可以转换为其他素数计数函数的显式表达式。

根据前节（第~\ref{sec:integration-by-parts} 节）的推导，$J$ 的显式公式为（公式 \eqref{eq:explicit_J_final}；
$x$ 非素数幂时 $J_0=J$）：
\[
J_0(x)=\li(x)-\sum_{\rho}\li(x^{\rho})-\log 2+\int_x^{\infty}\frac{dt}{t(t^2-1)\log t}.
\]
以此为基础，本节进一步通过 Möbius 反演推导 $\pi$ 的显式公式（跳跃点用 $\pi_0$）。

\subsection{\texorpdfstring{$\pi(x)$}{pi(x)} 的显式公式}\label{sec:explicit-pi}

利用 Möbius 反演（第~\ref{ch:elementary_number_theoretic_functions}~章 \eqref{eq:pi_from_J}；跳跃点取平均时两端写 $\pi_0,J_0$）：
\[
\pi_0(x)=\sum_{n=1}^{N(x)}\frac{\mu(n)}{n}J_0(x^{1/n}),
\qquad N(x)=\left\lfloor\frac{\log x}{\log 2}\right\rfloor
\]
（$n>N(x)$ 时 $x^{1/n}<2$，$J_0=0$；与 \eqref{eq:pi_from_J} 同一截断）。

将 $J(x^{1/n})$ 的显式公式（公式 \eqref{eq:explicit_J_final}，以 $x^{1/n}$ 替代 $x$）代入每一项：
\[
J(x^{1/n}) = \li(x^{1/n}) - \sum_\rho \li(x^{\rho/n}) - \log 2
              + \int_{x^{1/n}}^\infty \frac{dt}{t(t^2-1)\log t}.
\]

对 $n=1,2,\ldots,N(x)$ 逐项代入后求和。


\paragraph{级数交换次序的合法性}
当 $n$ 求和为有限和（$n \le N(x)$）时，与 $\rho$ 求和交换次序无条件成立。
因此 $\sum_{n=1}^{N(x)}\frac{\mu(n)}{n}\sum_\rho \li(x^{\rho/n})
= \sum_\rho \sum_{n=1}^{N(x)}\frac{\mu(n)}{n}\li(x^{\rho/n})$ 是合法的有限次序交换。
最终结果在每个固定的 $x$ 下均为有限求和，不涉及条件收敛的微妙问题。


对 $n=1,2,\ldots,N(x)$ 各项代入 \eqref{eq:explicit_J_final} 后加总，得有限展开
（跳跃点取左右平均时，左端写 $\pi_0(x)$；非素数处 $\pi_0=\pi$）：
\begin{equation}
\begin{aligned}
\pi_0(x) &= \sum_{n=1}^{N(x)}\frac{\mu(n)}{n}\,\li(x^{1/n})
         - \sum_{n=1}^{N(x)}\frac{\mu(n)}{n}\sum_{\rho}\li\!\left(x^{\rho/n}\right)\\
       &\quad{} - \log 2\cdot\sum_{n=1}^{N(x)}\frac{\mu(n)}{n}
         + \sum_{n=1}^{N(x)}\frac{\mu(n)}{n}\int_{x^{1/n}}^{\infty}\frac{dt}{t(t^2-1)\log t},
\end{aligned}
\label{eq:explicit_pi}
\end{equation}
其中 $N(x)=\lfloor\log x/\log 2\rfloor$，固定 $x$ 时仅有限项非零。

\paragraph{与 $R(x)$ 形的关系}
定义 $R(x):=\sum_{n=1}^{\infty}\frac{\mu(n)}{n}\li(x^{1/n})$（Gram）。
在合适整理下（含平凡零点贡献的闭式，Edwards §1.16），
\eqref{eq:explicit_pi} 与经典形
\[
\pi_0(x)=R(x)-\sum_{\rho}R(x^{\rho})-\frac{1}{\log x}+\frac{1}{\pi}\arctan\frac{\pi}{\log x}
\]
（von Mangoldt 等整理；原稿给出 $f$ 形与反演思想，后世写成 $R$ 形）相容。
$\sum_\rho$ 仍对称截断、计重数。


\section{显式公式的物理诠释：零点振荡}\label{sec:physical-interpretation}

显式公式的结构特征如下：
\[
\psi_0(x) = \underbrace{x}_{\text{主项}} - \underbrace{\sum_{\rho}\frac{x^{\rho}}{\rho}}_{\text{零点振荡项}} - \underbrace{\log(2\pi) - \frac12\log(1-x^{-2})}_{\text{低阶修正}}
\]

\textbf{主项 $x$}：给出了素数定理的渐近主项 $\psi(x)\sim x$。

\textbf{零点振荡项}：$\rho=\beta+i\gamma$ 与 $\bar\rho$ 合并后，
\[
\frac{x^{\rho}}{\rho}+\frac{x^{\bar\rho}}{\bar\rho}
=\frac{2x^{\beta}}{|\rho|}\cos(\gamma\log x-\arg\rho)
\]
（$\beta\in(0,1)$ 时 $\arg\rho\in(-\pi/2,\pi/2)$）。
RH 下 $\beta=1/2$，振幅 $O(x^{1/2}/|\rho|)$；若某 $\beta>1/2$，振幅更大。
这是 RH 与素数误差阶的直接联系。细节见第~\ref{sec:convergence}~节。

\section{级数收敛性说明}\label{sec:convergence}

级数 $\sum_{\rho}x^{\rho}/\rho$ \textbf{不绝对收敛}，须对称截断：
\[
\sum_{\rho}\frac{x^{\rho}}{\rho}
=\lim_{T\to\infty}\sum_{0<\Im\rho<T}
\Biggl(\frac{x^{\rho}}{\rho}+\frac{x^{\bar\rho}}{\bar\rho}\Biggr)
\]
（计重数）。由 $N(T)\sim\frac{T}{2\pi}\log\frac{T}{2\pi e}$ 知 $\sum_{|\gamma|\le T}1/|\rho|$ 发散，
而共轭合并后的余弦因子使级数条件收敛。


\paragraph{Weierstrass-Hadamard 理论与逐项积分的严格性保障}
在黎曼最初的推导中，他是直接从 $\log \zeta(s)$ 的围道积分出发，通过将 $\log \zeta(s)$ 展开为关于非平凡零点的对数级数并进行“逐项积分”来获得显式公式的。这一操作在分析学上极具挑战性：因为在积分路径 $\Re(s)=a>1$ 上，当零点的模趋于无穷时，对数级数的各分项并不能保证一致收敛，且伴随复对数分支割线的困难。

正是 \textbf{Weierstrass 因子定理} 和 \textbf{Hadamard 因子分解定理} 提供了解决此困难的唯一合法桥梁：
\begin{itemize}
  \item Weierstrass 理论表明整函数可以通过其零点乘积结合指数修正因子来完整重构；
  \item Hadamard 定理则确立零点密度与\textbf{对称截断}下乘积/级数的收敛方式，将形式对数展开转化为局部一致收敛的解析级数；
  \item 由此，对数级数在有限截断路径 $[-iT, iT]$ 上的“逐项积分”获得了合法性，并且通过零点密度的先验估计，使得 $T \to \infty$ 时的极限换序最终能够被严格控制。
\end{itemize}
没有这套理论，显式公式中的逐项积分只能被视为一种启发式猜想；有了它的理论支撑，黎曼的“形式级数”才成为严格成立的解析等式。


\section{旁支：RH 与 Mertens 函数 \texorpdfstring{$M(x)$}{M(x)}}
\label{sec:mertens-rh}
\label{sec:mertens_rh}
\label{sec:mertens-equivalence}

设 $M(x)=\sum_{n\le x}\mu(n)$（Mertens 函数；定义只依赖 $\mu$，不依赖 $\zeta$ 的延拓）。
本节给出 RH 的一种\textbf{初等数论等价}（旁支：不进入 $\psi\to J\to\pi$ 主链，证明提要级）。

\begin{theorem}{Mertens 函数与黎曼猜想的等价性}{}
\label{thm:mertens_rh}
以下两命题等价：
\begin{enumerate}
  \item \textbf{RH}：全部非平凡零点满足 $\operatorname{Re}(\rho)=1/2$；
  \item \textbf{Mertens 界}：对任意 $\varepsilon>0$，$M(x)=O(x^{1/2+\varepsilon})$（$x\to\infty$）。
\end{enumerate}
\end{theorem}

\paragraph{Perron 表示}
由 $\sum\mu(n)n^{-s}=1/\zeta(s)$（$\operatorname{Re}s>1$）与 Perron 公式，
\begin{equation}
  M(x)
  \;=\;
  \frac{1}{2\pi i}\int_{c-i\infty}^{c+i\infty}
  \frac{1}{\zeta(s)}\,\frac{x^s}{s}\,\mathrm{d}s,
  \qquad c>1.
  \label{eq:Mx_perron}
\end{equation}
与 $\psi$ 不同：$1/\zeta$ 在 $s=1$ 无极点（故无主项 $x$），
奇点来自 $\zeta$ 的零点与因子 $1/s$ 在 $s=0$ 的极点（留数 $-2$）。
因而 $M(x)/x\to 0$ 与素数定理同阶事实相容，而 $M$ 的振荡由零点控制。

\paragraph{RH $\Rightarrow$ Mertens 界（提要）}
在 RH 下将~\eqref{eq:Mx_perron} 的路径左移至 $\operatorname{Re}s=1/2+\delta$（$\delta>0$ 任意小）：
该竖带内无非平凡零点，无额外留数。
在 $\operatorname{Re}s=1/2+\delta$ 上需 $|1/\zeta(s)|=O(|t|^{\varepsilon})$
（RH 下的经典界，Littlewood；见 Titchmarsh §14.3——\textbf{不能}由 $|\zeta|=O(|t|^A)$ 直接取倒数得到）。
路径积分 $\ll x^{1/2+\delta}T^{\varepsilon}$；取 $T=x$ 并令 $\delta,\varepsilon$ 充分小，
配合 Perron 截断误差 $O(x^{c}/T)$（第~\ref{ch:integral_transforms}~章），
得 $M(x)=O(x^{1/2+\varepsilon_0})$ 对任意 $\varepsilon_0>0$。
（更精细的 $O(x^{1/2}(\log x)^C)$ 亦属 Littlewood，此处不展开。）

\paragraph{Mertens 界 $\Rightarrow$ RH}
Mellin--Stieltjes 表示（$\operatorname{Re}s>1$）
\begin{equation}
  \frac{1}{\zeta(s)}
  \;=\;
  s\int_1^\infty\frac{M(x)}{x^{s+1}}\,\mathrm{d}x
  \label{eq:mellin_Mx}
\end{equation}
在 $M(x)=O(x^{1/2+\varepsilon})$ 时，右端积分对任意 $\delta>0$ 在 $\operatorname{Re}s\ge 1/2+\delta$ 绝对收敛，
故给出 $1/\zeta$ 在 $\operatorname{Re}s>1/2$ 上的全纯延拓。
从而 $\zeta$ 在 $\operatorname{Re}s>1/2$ 无零点，即全部非平凡零点满足 $\operatorname{Re}(\rho)\le 1/2$。
由 $\xi(s)=\xi(1-s)$（第~\ref{ch:xi_function}~章），零点关于 $\operatorname{Re}s=1/2$ 对称：
若某 $\rho$ 有 $\operatorname{Re}(\rho)<1/2$，则 $\operatorname{Re}(1-\rho)>1/2$，矛盾。
故 $\operatorname{Re}(\rho)=1/2$，即 RH。$\square$


%%%%%%%%%%%%%%%%%%%%%%%%%%%%%%%%%%%%%%%%%%%%%%%%%%%%%%%%%%%%
\section*{本章小结与后续展望}\label{sec:riemann_explicit_summary}
\addcontentsline{toc}{section}{本章小结与后续展望}

黎曼显式公式建立了 $\psi_0$（非素数幂处即 $\psi$）与非平凡零点之间的解析对应。
主项 $x$ 对应历史上 Legendre 型与 Gauss 对数积分等光滑逼近所抓住的大尺度趋势；
零点级数 $\sum x^\rho/\rho$ 给出此前连续近似所无法解释的振荡修正。
表示既已精确，仍须区分三层：公式给出形状；素数定理给出主项主导（$\psi\sim x$）；RH 约束振荡量级。
有公式不等于已证素数定理，亦不等于已证 RH。
素数定理是渐近式（如 $\psi\sim x$ 或 $\pi\sim\Li$）；完整显式公式是恒等式（截断后才是带余项的近似）。
由显式公式推出 $\psi\sim x$，还须将 $\sum_{\rho}x^{\rho}/\rho$（或 $\sum_{\rho}\li(x^{\rho})$ 的相应贡献）估计为 $o(x)$，并非写出公式即自动得 PNT；
PNT 的标准证明也不依赖显式公式，但在附加估计下显式公式可以蕴含 PNT。
RH 确保零点振荡振幅为 $O(x^{1/2})$ 阶，将 $\psi$ 的起伏约束在 $O(\sqrt{x}\log^2 x)$ 量级，
并经 $J\to\pi$ 得到 $\pi-\Li=O(\sqrt{x}\log x)$。

从初等算术层经解析延拓、$\xi$ 与零点理论到显式公式，\textbf{主体理论链条}至此完成；
全文着色的知识谱系总图见下一章（第~\ref{ch:spectrum_map}~章图~\ref{fig:function-relations}）。
（书名何以将「素数分布」与「黎曼猜想」并列，见书末后记。）

\vspace{1em}
\noindent\textbf{本章遗留的技术细节与参考文献：}
\begin{itemize}
  \item Perron 公式的截断误差估计：参见第\ref{ch:integral_transforms}章第\ref{subsec:perron_truncation}节。
  \item 无零点区域定理与 $\zeta'/\zeta$ 的界估计：参见 Titchmarsh (1986) §3.9, §9.4。
  \item $\li(x^\rho)$ 的复解析延拓与级数收敛性：参见 Edwards (1974) §1.14。
  \item 从 $\psi(x)$ 到 $J(x)$ 的常数合并完整计算：参见 Edwards (1974) §1.15。
\end{itemize}

下一章给出主体理论完成后的知识谱系全文着色总图；
其后各章再把复零点产生的高频振荡重组为对数余弦级数，并与经典 Fourier 逼近对比。

\begin{center}
\fbox{\parbox{0.92\textwidth}{
\centering
\vspace{0.2em}
\textbf{第~\ref{ch:riemann_explicit_formula}~章：重要定理与核心公式汇总}\\[6pt]
\renewcommand{\arraystretch}{1.6}
\begin{tabular}{lp{8.6cm}}
\hline
\textbf{名称} & \textbf{数学内容 / 公式表达} \\
\hline
$J(x)$ 黎曼显式公式 & $J(x) = \li(x) - \displaystyle\sum_\rho \li(x^\rho) - \log 2 + \displaystyle\int_x^\infty \dfrac{dt}{t(t^2-1)\log t}$ \\
$\psi(x)$ 显式公式 & $\psi(x) = x - \displaystyle\sum_\rho \dfrac{x^\rho}{\rho} - \log(2\pi) - \dfrac{1}{2}\log(1-x^{-2})$ \\
Mertens 函数 RH 等价定理 & $M(x) = \displaystyle\sum_{n \le x} \mu(n) = O(x^{1/2+\varepsilon}) \iff \operatorname{Re}(\rho) = 1/2$ \\
误差项 von Koch 界限 & $\pi(x) - \li(x) = O(\sqrt{x}\log x) \iff \operatorname{Re}(\rho) = 1/2$ \\
\hline
\end{tabular}
\vspace{0.2em}
}}
\end{center}